\scanpage{180}
\mspage{125}
% archive sheet number: 179

de b \unsure{(toutefois)|(resp.)}, b$'$, ou b$''$\deleted{\unsure{b|)}}, il reste à s'assurer que le foncteur $f^*_{W_S}$ décrit de $=$ le foncteur induit par le précédent sur les types d'hom. relatifs loc. \unsure{constants}. Or on a les inclusions canoniques $\unsure{\mathcal F|F}_0\subset\mathcal F$, $\mathcal F'_0\subset\mathcal F'$, et sur $\mathcal F_0,\mathcal F'_0$ $W_S=W_S^d\underset{\unsure{2}}=W_S^{\unsure{t|f}}$ et ces \deleted{\unsure{quatre}} catégories sont OK pour \deleted{calculer} décrire $\mathrm{HOT}^{\unsure{\mathrm{lc}|\mathrm{le}}}_{W}(S\deleted{\unsure{\,) }})$ et $\mathrm{HOT}^{\mathrm{lc}}_W(S')$, et le foncteur $\unsure{g^*|f^*}$ induit entre eux par images inverses.

% slanted
\margin{NB $(\mathrm{Cat}_{/S})_0=\underline{\unsure{\mathrm{Fibhot}|\mathrm{Fiblct}}}\,S$, pour les autres cas c'est clair.}

(\uline{NB} Le foncteur chgt. de base $X\mapsto X\times_S S'$ est compatible avec les localiseurs sur ces catégories, car c'est trivial sous la forme $W_S^f$.)

On a alors, par images inverses \inserted{$X\mapsto X\times_S S'$}, un diagramme commutatif

\diagram{\begin{tikzcd}[ampersand replacement=\&,column sep=large]
\mathcal F_0(W_S)^{-1}\arrow[r,"\approx"]\arrow[d,"{(f^*_{W_S})_0}"'] \& \mathcal F_{\deleted{\unsure{0}}}(W_S)^{-1}\arrow[d,"{f^*_{W_S}}"]\\
\mathcal F'_0(W_{S'})^{-1}\arrow[r,"\approx","{\deleted{\unsure{\simeq}}}"'] \& \mathcal F'(W_{S'})^{-1}
\end{tikzcd},}

où les flèches horizontales sont des \uline{équivalences}\unsure{\,/s} de catégories, par le \unsure{th. satellite} 1 \textcircled{I}. Cela prouve que $f^*_{W_S}$ lui aussi est bon pour calculer les images inverses de t. d'hom. relatifs.

\deleted{\unsure{On}} On a prouvé que \deleted{b$''$ \unsure{$\Longrightarrow$} b} b, b$'$, b$''$ impliquent chacun a), prouvons que a les implique. Or a) dit exactement que les images inverses \deleted{\unsure{par}} $g^*$ des éléments de $\mathrm{HOT}^{\mathrm{lc}}_W(S)$
